\section{附录：测度与统计方法速查}\label{app:metrics-stats-cheatsheet}

本附录以速查形式汇总本文的核心测度与统计学方法，目的为：在不改变主线研究问题（RQ1--RQ3）主终点定义的前提下，帮助审稿人快速理解「每个指标测什么、在哪个口径有效、如何统计检验、哪些属于主分析/敏感性分析」。

\subsection{测度总表}

\small
\begin{longtable}{p{3cm} p{2.5cm} p{2cm} p{1.5cm} p{3cm} p{3.5cm}}
\caption{测度总表（符号、层级、适用口径与用途）}\label{tab:metrics-overview}\\
\toprule
\textbf{测度} & \textbf{符号/字段} & \textbf{层级} & \textbf{口径} & \textbf{主要用途} & \textbf{说明/边界}
\\
\midrule
\endfirsthead

\multicolumn{6}{l}{\small（续表~\ref{tab:metrics-overview}）}\\
\toprule
\textbf{测度} & \textbf{符号/字段} & \textbf{层级} & \textbf{口径} & \textbf{主要用途} & \textbf{说明/边界}
\\
\midrule
\endhead

\midrule
\multicolumn{6}{r}{\small 续下页}\\
\endfoot

\bottomrule
\endlastfoot

活跃标注时间 & \texttt{active\_time} & task / worker & T / I & RQ1 主终点（效率）；可作为“成本维度”辅助汇总 & 仅度量交互活跃时间；需配套审计表披露 unknown/缺失与多 session；若用于综合分需披露覆盖率与异常率。\\
改动幅度（工作量代理） & $\mathrm{IoU}_{\mathrm{edit}}(t,u)$ & annotation & T / I & RQ1 风险控制；反例 Type 2 & 仅在展示初始化（semi）任务定义；不表示正确性。\\
一致性（任务内） & $\mathrm{IAA}_t$ & task & I & RQ2 主终点（质量/可靠性变化） & 定义为任务内 pairwise IoU 的中位数；在预注册配对子集上做配对对照。\\
LOO 一致性代理 & $\mathrm{IoU}(A_{t,u},C_t^{(-u)})$ & annotation & I & RQ3 指标（离线回放） & 以 LOO 共识 $C_t^{(-u)}$ 为参考；要求 $|U_t|\ge 3$，更稳健时建议 $\ge 4$。\\
结构差异信号（角点对数差） & $\Delta n_{pairs}(t,u)$ & annotation / task & T / I & 分层报告；反例/失败溯源 & 仅作诊断与分层变量；不作为硬门控或主终点。\\
边界 RMSE（几何差异） & $\mathrm{BoundaryRMSE}$ & annotation & I / M & 质量/反例诊断（可比型门控后） & 用边界曲线差异替代不稳定点匹配；若曲线不可构造则记为门控失败并披露。\\
角点匹配 RMSE（辅助） & $\mathrm{RMSE}_{\mathrm{match}}(t,u)$ & annotation & I / M & 质量诊断（可比型门控） & 仅在匹配稳定且覆盖率足够时启用；否则记为 NA 并披露缺失率。\\
全局可靠度（工人） & $r_u$（median） & worker & I & 路由排序；画像纵轴 & 仅用 \texttt{Calibration\_manual} 估计；报告 95\% CI 与 LCB，用于保守排序。\\
场景特异可靠度 & $r_u^{(s)}$（median） & worker $\times$ scene & I & 核心场景路由（RQ3） & 仅在核心场景 $|\mathcal{S}_{\mathrm{core}}|\le 4$ 启用；不足门槛则退化为 $r_u$ 并披露启用/退化率。\\
场景启用率/退化率 & \texttt{activation\_rate} / \texttt{degeneration\_rate} & scene / worker & I & 回答“场景分层是否静默失效” & 作为审计性披露项：不足门槛或 CI 过宽导致退化到全局 $r_u$。\\
门控失败率 & \texttt{gate\_fail\_rate} & task / worker & T / I / M & 安全性/数据可用性审计（非主终点） & 区分可计算型失败（影响口径 M）与可比型失败（仅影响对应列）；用于证明不静默过滤与失败原因可追溯，而非作为路由收益主证据。\\
预算效率（序贯冗余） & $k_{used}(t)$ / $\bar{k}_{used}$ & task / split & T / I & RQ3 指标（预算效率） & 与 $k_0,k_{\max}$ 配套报告；需披露达到 $k_{\max}$ 的比例（专家复核/单独审计队列）。\\
标注前分布风险（OOD 代理） & $d_t$ / $I_t^{\mathrm{OOD}}$ & task & T / I & stress mode 触发；分层报告 & 主实现：共享 pre-head latent $\rightarrow \mathrm{GAP}_W \rightarrow$ L2 归一化 $\rightarrow$ 第 $K$ 近邻距离；$I_t^{\mathrm{OOD}}=\mathbf{1}[d_t>\tau_d]$，其中 $\tau_d$ 由校准集 leave-one-out 分位预注册锁定；不用于静默过滤。\\
标注前结构风险（诊断触发） & $g_t$ & task & T / I & stress mode 触发；失败溯源 & 由初始化结构合法性/失败模式构造；不可计算时记 NA 并披露。\\
工人风险层级 & $R_u^{\mathrm{tier}}$ & worker & T / I & 画像横轴；审计/降级规则 & 离散风险层级（Stable / Trust-vulnerable / Condition-fragile / Unusable-Noise）；不使用单一连续异常分数作为主画像轴。\\
风险签名 & $T_u,C_u,G_u$（主）; $E_u,M_u$（辅） & worker & T / I & 支撑画像分组与路由解释 & 主风险分别对应盲信初始化、条件性脆弱、可计算失败倾向；低努力与缺失仅作辅助审计。\\
标准 Layout 指标（对齐基准） & 2D IoU / 3D IoU & task & M & 与 HoHoNet 等对齐的补充报告 & 仅在口径 M（可计算型门控后）报告；不用于样本静默剔除。\\
渲染深度指标（对齐基准） & depth RMSE / $\delta$ & task & M & 补充质量刻画（几何一致性） & 基于 layout 渲染得到深度图再比较；若渲染失败记为门控失败并披露。\\
OOS 比例与子类 & \texttt{scope=OOS} + 子类 & task & T & 数据可用性审计 & OOS 不进入主指标（I/M），但必须在 T 中报告比例、子类分布与冲突率。\\
多选标签共识/分歧 & $c_{t,k}$, $\mathrm{conf}_{t,k}$, $\mathrm{disagree}_t$ & task $\times$ tag & T / I & 场景代理；序贯追加诊断 & 仅用于分层、审计与追加触发（任务已获得足够元标签后）；不得用于静默过滤样本。\\
一致性统计（标签） & Fleiss' $\kappa$（per-tag） & tag / split & T / I & 证明 meta-label 可解释性 & 低频标签需单独标注 low-freq 警告；仅作审计性解释而非硬门控。\\

\end{longtable}
\FloatBarrier

\subsection{统计方法速查（主终点 + 辅助检验 + 敏感性分析）}

\begin{table}[H]
\centering
\caption{统计方法速查表（对应 RQ、主终点、检验与 CI）}
\label{tab:stats-overview}
\small
\begin{tabular}{p{0.10\linewidth} p{0.24\linewidth} p{0.20\linewidth} p{0.22\linewidth} p{0.20\linewidth}}
\toprule
\textbf{RQ} & \textbf{主终点/对比} & \textbf{效应量（主）} & \textbf{检验（主）} & \textbf{区间与备注}
\\
\midrule
RQ1 & Manual vs Semi 的 \texttt{active\_time} & Hodges--Lehmann 中位数差 或相对改进 & Mann--Whitney $U$（非参数） & 95\% BCa bootstrap CI（重采样次数预注册）；同时报告与 $\mathrm{IoU}_{\mathrm{edit}}$ 的关系作为风险控制。\\
RQ2 & 配对子集（$N_{img}=25$）的 $\mathrm{IAA}_t$：manual vs semi & $\mathrm{median}(\Delta_t)$，$\Delta_t=\mathrm{IAA}_t^{semi}-\mathrm{IAA}_t^{manual}$ & 同图配对置换检验（随机翻转符号） & 配对 bootstrap 95\% CI；K--S 仅附录敏感性分析，不作为主终点。\\
RQ2(反例) & 同配对子集的反例类型分布差异 & 类别比例差 / OR & $\chi^2$ 检验（或小样本用 Fisher） & 报告效应量与 CI（必要时 bootstrap）；与主终点一并计入多重检验校正。\\
RQ2(加权共识) & 加权 vs 未加权 $\kappa$/accuracy & $\Delta\kappa$ / $\Delta\mathrm{acc}$ & BCa bootstrap 重采样（任务级） & CI 不跨 0 且效应量 $\geq MDE_{\kappa}$ 认为显著；必要时用 McNemar 比较反例检出差异。\\
RQ3(离线回放) & 策略 3 vs 策略 1/2 的一致性代理、预算效率 & 任务级差异的中位数/均值差；$\bar{k}_{used}$ 差异 & 主分析优先配对置换/配对 bootstrap（t-test 仅附录敏感性） & 门控失败率与失败原因分布作为安全性审计披露（给出比例与 CI），不纳入主检验集合；若部署仅执行策略 3，则策略 1/2 仅在支持集做影子评估并披露支持率。\\
RQ3(部署) & $Validation_{semi}$ 上策略 3 的一致性代理（若可）、反例拦截、预算效率 & 差异/比值 & 配对置换/配对 bootstrap；二元率以 binomial CI 描述性披露 & 门控失败作为安全性事件仅报告比例、原因分布与上界风险，不作为“收益显著性”主张来源。\\
多重检验 & 计划检验集合（预注册） & --- & Bonferroni（保守） & 只对预注册检验集合校正；附录探索性分析单独标注，不与主结论混用。\\
\bottomrule
\end{tabular}
\end{table}

\subsection{方法实现细节速查（举例说明实现步骤）}

\begin{itemize}
	\item \textbf{加权共识的预注册实施：} 以 \texttt{PreScreen\_manual} 中 20--22 张 manual anchors 作为独立锚点，按第~\ref{sec:method-weighted-consensus} 节的方案 A（重复平衡二折交叉拟合作为主锁定方案）在每折内估计 $r_u^{(0)}$ 与初始 $w_u$，并在验证折上比较加权与未加权一致性（$\Delta\kappa$ 或 $\Delta\mathrm{acc}$）。3 折分层切分仅作附录敏感性分析；主分析不采用 5 折。选出跨重复切分均值最优且满足 $w_{max}\in\mathcal{W}$ 的值后即冻结，不再事后调参。若样本不足，使用方案 B 的 A/B 拆分方案，拆分种子与结果明细都需预注册并公开。
	\item \textbf{可靠度与场景特异性的 LOO 估计：} $r_u$ 与 $r_u^{(s)}$ 全部基于 \texttt{Calibration\_manual} 的 $|U_t|\geq3$ 样本，计算 LOO 共识 $C_t^{(-u)}$ 避免自我影响；若某场景启用次数不足（$|\mathcal{S}_{\mathrm{core}}|\leq4$），退化为全局 $r_u$ 并在附录中报告启用/退化率。置信区间由 bootstrap 获得，若半宽未达设定精度（预注册置信宽阈值），即时从 \texttt{Calibration\_reserve} 中追加任务。
	\item \textbf{序贯停止与预算效率计算：} Validation 中每张任务在候选池内按 $d_t$/$g_t$ 分层并触发应激模式，序贯增加 $k$ 直到 \texttt{stop\_threshold} 或 $k_{\max}$。平均每任务使用的 $\bar{k}_{used}$ 通过 paired bootstrap 比较；门控失败仅作为安全性审计事件记录其比例与原因分布，不纳入主检验集合。
	\item \textbf{Kolmogorov--Smirnov (K--S) 两样本检验（敏感性分析）：} 用于比较两组样本的分布（例如 $\mathrm{IAA}_t$ 在 manual vs semi 的分布形状）。计算经验分布函数 $F_n,G_m$ 并取最大差值 $D_{n,m}=\sup_x|F_n(x)-G_m(x)|$。实现上优先使用 \texttt{scipy.stats.ks\_2samp}（two-sided）；当存在大量并列值或样本量较小时，改用置换检验：在原样本上随机交换组标签 $B_{perm}$ 次（预注册建议 $B_{perm}=2000$ 或 $10000$）构造置换分布并计算双侧 p-value。报告原始 $D$ 值、p-value 与置换分布摘要，并在附录注明 K--S 对离散/并列数据的局限性。
	\item \textbf{Cohen's d（配对样本）：} 对配对连续量（如同任务下的均值差）计算配对效应量。令差分 $D_i=X_i^{(A)}-X_i^{(B)}$，配对 Cohen's d 定义为 $d_{paired}=\bar{D}/s_D$，其中 $s_D$ 为差分的样本标准差。实现要点：同时报告 $d$ 的 95\% CI，使用 paired bootstrap（预注册重采样次数 $B_{boot}=5000$ 或 10000）获得 CI；当差分分布严重偏离正态时，补报 Hodges--Lehmann 中位数差与其 bootstrap CI 作为稳健替代。解释阈值参考 small=0.2, medium=0.5, large=0.8，但以 MDE 与研究语境为准。
	\item \textbf{多重检验与 MDE 门槛：} 7 个预注册检验在 Bonferroni 下校正（$\alpha\approx 0.0071$），仅在 $p<\alpha$ 且效应量超过预注册 MDE（例如 Time $\geq8\%$ 改进、一致性差 $\geq0.05$、以及预算效率的最小可解释改进）时宣布显著；门控失败率不进入主检验集合，仅作安全性/可用性审计披露。MDE 取自 Pilot → PreScreen 的分级锁定流程，整个过程中不得事后调整阈值。
	\item \textbf{指标数据采集与一致性检验：} \texttt{active\_time} 以 session 分段追踪（见第~\ref{sec:method-active-time} 节），每种汇总（per session/per annotator/per task）都按最长累计值汇总，并通过 Spearman 相关与原生 \texttt{lead\_time} 比值的中位数、IQR 验证采集逻辑一致性。
\end{itemize}

\subsection{关于加权几何平均综合分的预注册使用边界}

若需报告综合分以帮助读者快速把握整体趋势（类比综合跑分），建议将其定位为\textbf{辅助汇总指标}而非替代 RQ1--RQ3 主终点。为降低调参空间，权重与归一化方式需在实验前冻结，且不得用于选择样本、选择策略或调参。

\paragraph{审稿友好的“最小自由度”建议（权重与敏感性分析）}
\textbf{(1) 只对同一研究问题内的输出做汇总：}优先将综合分限定在 RQ3（路由策略对比）的输出维度内（例如一致性代理与预算效率；若该集合确有可靠的时间日志，可额外加入时间维度）。门控失败率作为安全性/可用性审计项单独披露，不建议纳入综合分（稀有事件在小样本下会导致不稳定权重争议）。不建议跨 RQ1--RQ3 把时间、质量、路由收益混成单一总分，避免因口径与缺失机制不同引发选择偏差争议。

	\textbf{(2) 权重只给少量离散方案并全部报告：}避免连续权重“网格搜索最优”。建议预注册并固定报告少量离散权重（示例，按 2 维：一致性代理/预算效率）：
\begin{itemize}
	\item 等权：$w=(0.50,\,0.50)$（默认；最不易被指控偏好性）。
	\item 质量优先：$w=(0.67,\,0.33)$（强调一致性代理）。
	\item 成本优先：$w=(0.33,\,0.67)$（强调预算效率）。
\end{itemize}
\textbf{若将 \texttt{active\_time} 纳入跑分：}建议仅在“用于综合分的任务集合”上满足以下条件时启用：时间字段覆盖率 $\ge 90\%$、unknown/缺失异常率低于预注册阈值、且通过 \texttt{active\_time} 与 \texttt{lead\_time} 的一致性审计（见主文 RQ1 的一致性校验）。此时将时间作为第 4 个“成本维度”（越小越好，进入综合分前需做单调变换为越大越好），并同样仅报告离散的 3 套权重，例如：
\begin{itemize}
	\item 等权3维：$w=(1/3,\,1/3,\,1/3)$（质量/预算/时间）。
	\item 质量优先3维：$w=(0.50,\,0.25,\,0.25)$。
	\item 成本优先3维：$w=(0.25,\,0.50,\,0.25)$。
\end{itemize}
其中时间维度的单调变换需预注册锁死（例如基于 baseline 的相对改进并做截断），并强制报告时间维度的覆盖率与异常率，避免因缺失机制引发样本切换。

\textbf{(3) 只做低成本稳健性披露：}附录报告 leave-one-dimension-out（每次去掉一个维度重算综合分）以及几何平均 vs 算术平均的替代聚合，检查结论方向与排序是否一致；并强制披露综合分的有效样本覆盖率（complete-case vs partial-case），避免因门控/缺失导致的隐性样本切换。